\begin{abstract}
A network intrusion detection system (NIDS) classifies heterogeneous network flows across severely imbalanced traffic categories. High overall accuracy on skewed traffic often conceals poor recognition of rare cyberattacks. The present study investigated how a standalone decision tree compared with decision-tree ensemble methods on a ten-class subset of the University of New South Wales network-based 2015 (UNSW-NB15) dataset comprising 257,673 cleaned observations and 41 predictors. An extended iterative dichotomiser three (ID3) learner served as the common base classifier across all models. The controlled design isolated the effects of post-pruning, bootstrap aggregation, and stagewise additive modelling using a multi-class exponential loss function (SAMME) from differences between unrelated model families.

The dataset was divided into a fixed stratified four-way split consisting of 123,683 growth rows, 41,225 reduced-error pruning rows, 41,234 validation rows, and 51,531 reserved generalisation rows. Preliminary learning-curve experiments up to 50,000 rows and a feature-correlation ablation showed that training sample size and minority class support strongly governed predictive performance. Removing strongly correlated predictors degraded validation macro-averaged F1-score.

A final comparison used all 123,683 growth rows and a shared hybrid class-balancing policy. On the reserved generalisation set, a single reduced-error-pruned decision tree achieved the strongest performance with a mean macro-averaged F1-score of $0.5993$, a balanced accuracy of $63.02\%$, and an overall accuracy of $80.73\%$. Bagging with 25 unpruned members improved over the standalone unpruned tree from $0.5545$ to $0.5747$ in mean macro-averaged F1-score. Boosting with 25 depth-eight SAMME members achieved a mean macro-averaged F1-score of $0.5408$. A supplementary validation extension showed no additional macro-averaged F1-score gain when expanding bagging from 25 to 50 members. The results demonstrated that combining hybrid class-balanced training with reduced-error post-pruning enabled a single decision tree to outperform the tested unpruned bagging and depth-bounded boosting ensembles.
\end{abstract}
